\section{Resultados e Discussão}

Nesta seção serão apresentados os resultados e será realizada uma análise comparativa dos modelos resultantes do
\textit{fine-tuning} realizado nas duas etapas de desenvolvimento. Além disso serão feitas comparações com o modelo
base e entre os métodos de treinamento.

\subsection{Desempenho na Classificação de Lesões de Pele}

A \autoref{tab:overall_comparison_skin_lesions} apresenta uma comparação entre todos os testes realizados.

\clearpage

\begin{table}[ht]
    \centering
    \begin{tabular}{l|l|c|c}
        \hline
        Tipo de imagem                  & Método de \textit{fine-tuning} & Volume de treinamento & Acurácia (\%) \\
        \hline
        \multirow{10}{*}{Dermatoscopia} & Nenhum (Quantizado)            & -                     & 11,8          \\
                                        & Nenhum                         & -                     & 14,0          \\
        \cline{2--2}
                                        & \multirow{3}{*}{\ac{qlora}}    & 450                   & 65,5          \\
                                        &                                & 900                   & 65,2          \\
                                        &                                & 5000                  & 85,4          \\
        \cline{2--2}
                                        & \multirow{3}{*}{\ac{lora}}     & 450                   & 49,5          \\
                                        &                                & 900                   & 42,0          \\
                                        &                                & 5000                  & \textbf{86,2} \\
        \hline
        \multirow{4}{*}{Aproximação}    & Nenhum (Quantizado)            & -                     & 14,8          \\
                                        & Nenhum                         & -                     & 17,0          \\
        \cline{2--2}
                                        & \ac{qlora}                     & 4736                  & \textbf{45,2} \\
        \cline{2--2}
                                        & \ac{lora}                      & 4736                  & 44,5          \\
        \hline
    \end{tabular}
    \label{tab:overall_comparison_skin_lesions}
    \caption{Acurácia na classificação de lesões de pele nos testes realizados. Os valores em
    negrito indicam as maiores acurácias para o tipo de imagem especificado. Fonte: Autoria própria.}
\end{table}

Os resultados obtidos na etapa final de desenvolvimento indicam que os modelos que passaram pelo \textit{fine-tuning}
para a classificação de lesões de pele com imagens de aproximação obtiveram um desempenho mediano, ficando abaixo do
esperado. Em comparação com os modelos da primeira etapa, treinados para a classificação com imagens de dermatoscopia,
a acurácia foi consideravelmente menor.

Os modelos da primeira etapa obtiveram um desempenho superior com um quantidade muito menor de dados, comparativamente,
o modelo treinado com \ac{qlora} e 5000 imagens de dermatoscopia superou o modelo treinado com \ac{qlora} e 4736
imagens de aproximação em 42,2\%. Embora existam diferenças nos hiperparâmetros de treinamento desses modelos, a
divergência é significativa demais para estar relacionada apenas com os hiperparâmetros.

\subsection{Desempenho na Classificação de Risco}

A classificação de risco introduzida na etapa final do trabalho apresentou um cenário mais promissor no desempenho,
porém a acurácia ainda foi mediana e pouco satisfatória. A \autoref{tab:overall_comparison_risk} apresenta uma visão
comparativa da acurácia na classificação das categorias de risco, evidenciando as melhorias no desempenho
proporcionadas pelo \textit{fine-tuning}.

\begin{table}[ht]
    \centering
    \begin{tabular}{l|c}
        \hline
        Método de \textit{fine-tuning} & Acurácia (\%) \\ \hline
        Nenhum (Quantizado)            & 23,4          \\
        Nenhum                         & 27,4          \\
        \ac{qlora}                     & 54,6          \\
        \ac{lora}                      & \textbf{55,1} \\ \hline
    \end{tabular}
    \label{tab:overall_comparison_risk}
    \caption{Acurácia na classificação de risco para os testes da etapa final de desenvolvimento. O valor em
    negrito indica a maior acurácia. Fonte: Autoria própria.}
\end{table}

\subsection{Desempenho na Geração de Laudos}

O desempenho na geração de texto foi considerado satisfatório, conforme pode ser viso na
\autoref{tab:overall_nlp_metrics}, tanto o modelo treinado com o \ac{qlora} quanto o com \ac{lora} atingiram um
desempenho superior ao do modelo base, constatando que os modelos conseguiram generalizar o formato de geração do texto
do laudo. Nas métricas \ac{bleu} e \ac{meteor} foi observado o maior ganho de desempenho em relação ao modelo base,
enquanto na métrica \ac{rouge} o ganho de desempenho foi menor.

\begin{table}[ht]
    \centering
    \begin{tabular}{l|c|c|c|c}
        \hline
                                       & \multicolumn{4}{c}{Método de \textit{fine-tuning}}
        \\ \hline
        Métrica                        & Nenhum (Quantizado)
                                       & Nenhum                                             & \ac{qlora}
                                       & \ac{lora}
        \\ \hline
        \ac{bleu}-1                    & 0,181                                              & 0,192      &
        \textbf{0,551}
                                       & 0,545
        \\
        \ac{bleu}-2                    & 0,125                                              & 0,129      &
        \textbf{0,524}
                                       & 0,517
        \\
        \ac{bleu}-3                    & 0,099                                              & 0,102      &
        \textbf{0,511}
                                       & 0,504
        \\
        \ac{bleu}-4                    & 0,081                                              & 0,083      &
        \textbf{0,500}
                                       & 0,493
        \\
        \hline
        \ac{rouge}-1 (\textit{recall}) & 0,422                                              & 0,417      &
        \textbf{0,678}
                                       & 0,677
        \\
        \ac{rouge}-L (\textit{recall}) & 0,330                                              & 0,313      &
        \textbf{0,629}
                                       & 0,627
        \\
        \hline
        \ac{meteor}                    & 0,211                                              & 0,221      &
        \textbf{0,580}
                                       & 0,576
        \\
        \hline
    \end{tabular}
    \label{tab:overall_nlp_metrics}
    \caption{Métricas de \ac{pln} do texto das respostas geradas sobre os testes na etapa final de
    desenvolvimento. Os maiores valores de cada métrica estão destacados em negrito. Fonte: Autoria própria.}
\end{table}

A fim de realizar uma análise aprofundada da geração de texto, as respostas abaixo foram coletadas dos resultados dos
testes. Primeiramente, o texto de referência é apresentado:

\begin{dialogue}
    \speak{Modelo} Classificação: Câncer de pele Não Melanoma (CEC ou CBC). \\
    Classificação de risco: AMARELA - ENCAMINHAMENTO COM PRIORIDADE PARA O AMBULATÓRIO DE REFERÊNCIA TERCIÁRIO. \\
    Conclusão sobre a lesão: compatível com câncer da pele. \\
    Conclusão: Encaminhar o paciente para o ambulatório de Dermatologia - CÂNCER DA PELE/Cirurgia Dermatológica ou
    referência regional. \\ \\
    Se o paciente já realizou biópsia da lesão, o resultado deve ser levado no dia da cirurgia.
\end{dialogue}

As duas respostas abaixo foram geradas pelo \ac{llama}-3.2-11B quantizado e não quantizado, respectivamente. Os modelos
classificaram corretamente a lesão de pele, mas apenas a versão quantizada categorizou o risco de forma correta. Além
disso, as conclusões não apresentaram as recomendações ideais.

\begin{dialogue}
    \speak{Modelo} Classificação: Câncer de pele Não Melanoma (CEC ou CBC). \\
    Classificação de risco: AMARELA - ENCAMINHAMENTO COM PRIORIDADE PARA O AMBULATÓRIO DE REFERÊNCIA TERCIÁRIO. \\
    Conclusão sobre a lesão: A lesão apresentada na imagem é caracterizada por uma lesão de pele com aspecto ulcerado e
    com bordas irregulares, o que sugere uma lesão maligna. A classificação de risco como amarela indica que a lesão
    requer
    acompanhamento e tratamento em um ambulatório de referência terciário. \\
    Conclusão: A lesão apresentada na imagem é suspeita de câncer de pele não melanoma e requer acompanhamento e
    tratamento em um ambulatório de referência terciário.
\end{dialogue}

\begin{dialogue}
    \speak{Modelo} Classificação: Câncer de pele Não Melanoma (CEC ou CBC). \\
    Classificação de risco: BRANCA - SEM NECESSIDADE DE INTERVENÇÃO OU ACOMPANHAMENTO. \\
    Conclusão sobre a lesão: A lesão apresentada na imagem é caracterizada por uma lesão de pele com uma borda
    irregular e um centro escuro, o que sugere uma lesão maligna. A presença de uma lesão de pele com essas
    características
    é compatível com um Câncer de pele Não Melanoma (CEC ou CBC). \\
    Conclusão: A lesão apresentada na imagem é compatível com um Câncer de pele Não Melanoma (CEC ou CBC) e deve ser
    avaliada por um especialista em dermatologia.
\end{dialogue}

As respostas seguintes foram geradas com o \ac{llama}-3.2-11B treinado com \ac{qlora} e \ac{lora}, respectivamente.
Ambos classificaram as lesões e as categorias de risco corretamente. As conclusões geradas seguem o texto de referência
com precisão, mas apresentam a palavra ``compatível'' com a grafia incorreta, omitindo o acento agudo. O texto gerado
também inclui um espaçamento inesperado.

\begin{dialogue}
    \speak{Modelo} Classificação: Câncer de pele Não Melanoma (CEC ou CBC). \\
    Classificação de risco: AMARELA - ENCAMINHAMENTO COM PRIORIDADE PARA O AMBULATÓRIO DE REFERÊNCIA TERCIÁRIO. \\
    Conclusão sobre a lesão: compativel com câncer da pele. \\
    Conclusão: Encaminhar o paciente para o ambulatório de Dermatologia - CÂNCER DA PELE/Cirurgia Dermatológica ou
    referência regional. \\ \\ \\ \\
    Se o paciente já realizou biópsia da lesão, o resultado deve ser levado no dia da cirurgia.
\end{dialogue}

\begin{dialogue}
    \speak{Modelo} Classificação: Câncer de pele Não Melanoma (CEC ou CBC). \\
    Classificação de risco: AMARELA - ENCAMINHAMENTO COM PRIORIDADE PARA O AMBULATÓRIO DE REFERÊNCIA TERCIÁRIO. \\
    Conclusão sobre a lesão: compativel com câncer da pele. \\
    Conclusão: Encaminhar o paciente para o ambulatório de Dermatologia - CÂNCER DA PELE/Cirurgia Dermatológica ou
    referência regional. \\ \\ \\ \\
    Se o paciente já realizou biópsia da lesão, o resultado deve ser levado no dia da cirurgia.
\end{dialogue}

\subsection{Comparação entre as Técnicas de Treinamento}

Por fim, comparando as métricas de treinamento apresentadas na \autoref{fig:trainig_time}, \autoref{fig:used_memory} e
\autoref{tab:final_training_metrics}, observa-se que ambos os métodos de treinamento tomam uma quantidade similar de
tempo para serem concluídos. O tempo de treinamento também tende a escalar na mesma taxa para ambos os métodos, mesmo
com diferentes volumes de dados. Ademais, o método \ac{lora} tente a utilizar uma quantidade maior de memória, com um
uso médio de 84,19\% a mais de memória em relação ao \ac{qlora} na etapa inicial. Na etapa final, o \ac{lora} utilizou
25,94\% a mais de memória. A disparidade entre essas diferenças para as duas etapas pode ser explicada pelas alterações
nos hiperparâmetros que afetam o uso da memória, como tamanho do lote por dispositivo e o passo de acumulação de
gradiente. Além disso, é possível que as características dos conjunto de dados tenham influenciado nesta diferença.

\begin{figure}[ht]
    \centering
    \includegraphics[width=0.8\columnwidth,keepaspectratio]{images/training_time.png}
    \label{fig:trainig_time}
    \caption{Tempo de treinamento dos modelos da etapa inicial com diferentes técnicas de treinamento. Fonte: Autoria
    própria.}
\end{figure}

\clearpage

\begin{figure}[ht]
    \centering
    \includegraphics[width=0.8\columnwidth,keepaspectratio]{images/used_memory.png}
    \label{fig:used_memory}
    \caption{Memória utilizada no treinamento dos modelos da etapa inicial com diferentes técnicas de
    treinamento. Fonte: Autoria própria.}
\end{figure}

\begin{table}[ht]
    \centering
    \begin{tabular}{l|c|c}
        \hline
        Método de treinamento & Tempo de treinamento (min) & Memória utilizada (GB) \\ \hline
        \ac{qlora}            & 286,83                     & 47,57                  \\
        \ac{lora}             & 276,46                     & 59,91                  \\ \hline
    \end{tabular}
    \label{tab:final_training_metrics}
    \caption{Métricas de treinamento para a etapa final do projeto. Fonte: Autoria própria.}
\end{table}
